% Part: incompleteness
% Chapter: incompleteness-provability
% Section: second-incompleteness-thm

\documentclass[../../../include/open-logic-section]{subfiles}

\begin{document}

\olfileid[pt]{inc}{inp}{2in}

\olsection{O segundo teorema da incompletude}

Como podemos expressar a afirmação de que $\Th{PA}$ não prova a própria
consistência? Dizer que $\Th{PA}$ é inconsistente equivale a dizer que
$\Th{PA} \Proves \eq[0][1]$. Portanto, podemos tomar a afirmação de
consistência $\OCon[\Th{PA}]$ como a !!{sentence} $\lnot
\OProv[\Th{PA}](\gn{\eq[0][1]})$; então o teorema seguinte cumpre o papel:

\begin{thm}
\ollabel{thm:second-incompleteness} 
Supondo que $\Th{PA}$ seja consistente, $\Th{PA}$ não !!{derive}
$\OCon[\Th{PA}]$.
\end{thm}

É importante observar que o teorema depende da representação específica de
$\OCon[\Th{PA}]$ (O.s.\@, da representação específica de
$\OProv[\Th{PA}](y)$). Usaremos apenas o fato de que a representação de
$\OProv[\Th{PA}](y)$ satisfaz as três condições de !!{derivability}; assim,
o teorema se generaliza para qualquer teoria com !!a{derivability}
predicado que tenha essas propriedades.

É esclarecedor ler o esboço do argumento de G\"odel, pois o teorema surge
como uma boa conclusão de efeito. O argumento é o seguinte. Seja
$!G_\Th{PA}$ a sentença de G\"odel que construímos na prova de
\olref[1in]{thm:first-incompleteness}. Mostramos: ``Se $\Th{PA}$ é
consistente, então $\Th{PA}$ não !!{derive} $!G_\Th{PA}$.'' Se
formalizarmos isso \emph{em} $\Th{PA}$, teremos uma prova de
\[
\OCon[\Th{PA}] \lif \lnot \OProv[\Th{PA}](\gn{!G_\Th{PA}}).
\]
Suponha agora que $\Th{PA}$ !!{derive}s $\OCon[\Th{PA}]$. Então ela
!!{derive}s $\lnot \Prov[\Th{PA}](\gn{!G_\Th{PA}})$. Mas, como
$!G_\Th{PA}$ é uma sentença de G\"odel, isso é equivalente a
$!G_\Th{PA}$. Portanto, $\Th{PA}$ !!{derive}s $!G_\Th{PA}$.

Mas sabemos que, se $\Th{PA}$ é consistente, ela não !!{derive}
$!G_\Th{PA}$!{} Assim, se $\Th{PA}$ é consistente, ela não pode !!{derive}
$\OCon[\Th{PA}]$.

Para tornar o argumento mais preciso, tomaremos $!G_\Th{PA}$ como a
sentença de G\"odel para~$\Th{PA}$ e usaremos as condições de
!!{derivability} (P1)--(P3) para mostrar que $\Th{PA}$ !!{derive}s
$\OCon[\Th{PA}] \lif !G_\Th{PA}$. Isso mostrará que $\Th{PA}$ não
!!{derive} $\OCon[\Th{PA}]$. Eis um esboço da prova, em~$\Th{PA}$.
(Para simplificar, omitimos os índices $\Th{PA}$.)
\begin{align}
& !G \liff \lnot \OProv(\gn{!G}) \ollabel{G2-1}\\
& \qquad\text{$!G$ é uma sentença de G\"odel}\notag \\
& !G \lif \lnot \OProv(\gn{!G}) \ollabel{G2-2}\\
  & \qquad\text{de \olref{G2-1}} \notag\\
& !G \lif
  (\OProv(\gn{!G}) \lif \lfalse) \ollabel{G2-3}\\
  & \qquad\text{de \olref{G2-2}, pela lógica}\notag\\
& \OProv(\gn{
    !G \lif
    (\OProv(\gn{!G}) \lif \lfalse)
  }) \ollabel{G2-4}\\
  & \qquad\text{de \olref{G2-3}, pela condição P1} \notag\\
& \OProv(\gn{!G}) \lif
  \OProv(\gn{
    (\OProv(\gn{!G}) \lif \lfalse)
    }) \ollabel{G2-5}\\
  & \qquad\text{de \olref{G2-4}, pela condição P2} \notag\\
& \OProv(\gn{!G}) \lif (\OProv(\gn{\OProv(\gn{!G})}) \lif \OProv(\gn{\lfalse})) \ollabel{G2-6}\\
  & \qquad\text{de \olref{G2-5}, pela condição P2 e pela lógica} \notag\\
& \OProv(\gn{!G}) \lif 
  \OProv(\gn{\OProv(\gn{!G})}) \ollabel{G2-7}\\
   & \qquad\text{por P3} \notag\\
& \OProv(\gn{!G}) \lif \OProv(\gn{\lfalse}) \ollabel{G2-8}\\
  & \qquad \text{de \olref{G2-6} e \olref{G2-7}, pela lógica}\notag\\
& \OCon \lif \lnot \OProv(\gn{!G}) \ollabel{G2-9}\\
  & \qquad\text{contraposição de \olref{G2-8} e $\OCon \ident \lnot \OProv(\gn{\lfalse})$}\notag \\
& \OCon \lif !G \notag\\
  & \qquad\text{de \olref{G2-1} e \olref{G2-9}, pela lógica}\notag
\end{align}
O uso da lógica acima consiste apenas em fatos elementares da lógica
proposicional; p.ex., \olref{G2-3} usa $\Proves \lnot!A \liff (!A\lif
\lfalse)$, e \olref{G2-8} usa $!A \lif (!B \lif !C), !A \lif !B
\Proves !A \lif !C$. O uso da condição~P2 em \olref{G2-5} e
\olref{G2-6} se apoia em instâncias de~P2,
$\OProv(\gn{!A \lif !B}) \lif
(\OProv(\gn{!A}) \lif \OProv(\gn{!B}))$. Na primeira, $!A \ident !G$ e
$!B \ident \OProv(\gn{!G}) \lif \lfalse$; na segunda, $!A \ident
\OProv(\gn{G})$ e $!B \ident \lfalse$.

A versão mais abstrata do segundo teorema da incompletude é a seguinte:

\begin{thm}
\ollabel{thm:second-incompleteness-gen} Seja $\Th{T}$ qualquer teoria
consistente e !!{axiomatized} que estenda $\Th{Q}$, e seja
$\OProv[\Th{T}](y)$ qualquer fórmula que satisfaça as condições de
!!{derivability} P1--P3 para~$\Th{T}$. Então $\Th{T}$ não !!{derive}
~$\OCon[T]$.
\end{thm}

\begin{prob}
Mostre que $\Th{PA}$ !!{derive}s $!G_{\Th{PA}} \lif \OCon[\Th{PA}]$.
\end{prob}

\begin{digress}
A moral da história é que nenhuma teoria consistente ``razoável'' para a
matemática pode !!{derive} sua própria afirmação de consistência. Suponha que
$\Th{T}$ seja uma teoria da matemática que contenha $\Th{Q}$ e o raciocínio
``finitário'' de Hilbert (seja o que for isso). Então nem toda a teoria
$\Th{T}$ pode !!{derive} a afirmação de consistência de $\Th{T}$ e,
a fortiori, o fragmento finitário tampouco pode !!{derive} a
afirmação de consistência de~$\Th{T}$. Nesse sentido, não pode haver uma
prova finitária de consistência para ``toda a matemática''.

Há alguma margem de interpretação do termo ``finitário'', e G\"odel, no
artigo de 1931, admite a possibilidade de que algo que consideremos
``finitário'' esteja fora dos tipos de matemática que Hilbert queria
formalizar. Mas G\"odel estava sendo generoso; hoje é difícil imaginar como
poderíamos encontrar algo que possa razoavelmente ser chamado de finitário,
mas que não seja formalizável, digamos, em $\Th{ZFC}$, a teoria dos conjuntos
de Zermelo--Fraenkel com o axioma da escolha.
\end{digress}

\end{document}
